%% ABGC: Bucket-level dye and Hodge projection clause selection
%%
%% Render with:
%%   lualatex -shell-escape abgc-mptp.tex && bibtex abgc-mptp && lualatex -shell-escape abgc-mptp.tex && lualatex -shell-escape abgc-mptp.tex

\documentclass[11pt]{article}

\usepackage{amsmath,amssymb,amsthm}
\usepackage{graphicx}
\usepackage[hyphens]{url}
\usepackage{hyperref}
\usepackage{booktabs}
\usepackage{longtable}
\usepackage{array}
\usepackage{tikz}
\usetikzlibrary{arrows.meta,positioning}
\usepackage{listings}
\lstset{basicstyle=\ttfamily\footnotesize,breaklines=true,frame=single,
        xleftmargin=1em,xrightmargin=0.5em}

\emergencystretch=3em

% Breakable monospace path / identifier
\newcommand{\lpath}[1]{\texttt{\nolinkurl{#1}}}
\newcommand{\lid}[1]{{\ttfamily\nolinkurl{#1}}}
\newcolumntype{L}[1]{>{\raggedright\arraybackslash}p{#1}}

\newtheorem{theorem}{Theorem}
\newtheorem{lemma}[theorem]{Lemma}
\newtheorem{definition}[theorem]{Definition}
\newtheorem{remark}[theorem]{Remark}
\newtheorem{observation}[theorem]{Observation}

\title{Bucket-level dye and Hodge projection clause selection:\\
       a +17 lift over \texttt{mzr02} on chainy MPTP,\\
       and what the HJB--NS refactor revealed}
\author{Mettapedia Working Notes}
\date{\today}

\begin{document}
\maketitle

\begin{abstract}
We report a clause-selection heuristic, \emph{ABGC} (Advective / Bidirectional /
Geodesic Control), which projects each E~prover unprocessed clause onto a
14-bit bucket signature, maintains a bidirectional dye field
$(a_\mathrm{fwd}, a_\mathrm{bwd})$ over buckets together with a directed flow
$u_{ij}$ across parent$\!\to\!$child bucket transitions, and applies a discrete
Hodge projection at batch boundaries before reading off a per-clause weight.
Interleaved at the 20\,\% slot of a 12-multiplier round-robin alongside
AI4REASON's \texttt{mzr02} Mizar-tuned strategy in E's HCB scheduler, the
controller solves $289/800$ on the validation slice of extended MPTP~5k
(\lpath{chainy_val_clean}) under a 30~s CPU limit per problem, four parallel
jobs, and a per-process 10~GiB virtual-memory cap. The matched
\texttt{mzr02}-alone single-heuristic baseline solves $272/800$ on the same
slice; the combined heuristic therefore lifts the baseline by $+17$ problems
($+6.2\,\%$ relative). E's portfolio \lid{--auto-schedule}, included for
context only, solves $334/800$.

A subsequent principled refactor mapped the controller onto the
Hamilton--Jacobi--Bellman / Navier--Stokes (HJB--NS) skeleton of
Goertzel and Stay: conjugate-gradient Hodge projection with isolated-node
handling, mass-conserving advection with a per-node CFL share, Laplacian
viscosity on the net-flow field, and graph-Laplacian diffusion of the dyes
with a degree-aware stability cap. The refactor was implemented end-to-end,
its conservation laws verified by telemetry, and a SEGV-class crash was
diagnosed and fixed (per-node CFL violation $\to$ negative dye $\to$
$\log(\mathrm{neg})$ $\to$ NaN in the clause heap $\to$ corrupted comparator
$\to$ SIGSEGV in \lid{ClauseSetExtractEntry}). The HJB--NS refactor did
\emph{not} reproduce the $+17$ lift: the best Sub-phase~A configuration
(conjugate-gradient projection without transport) achieves $253/800$, and
fully-wired Sub-phase~B transport achieves $246$--$238/800$, depending on
the projection sign convention.

We characterize the cause: the static-H Phase~4 hybrid implicitly reads a
\emph{divergence-amplified} signal in $u_{ij}$ that the clean Helmholtz
projection removes by design. The amplification was carried by the
under-converged Jacobi solver plus a sign convention on the projection
correction that adds rather than subtracts the gradient component. The
controller's H-term shape (\lid{meet}, \lid{align}, \lid{pressure}) was
implicitly tuned to that signal regime, and the principled transport
produces a structurally different signal that the current H-term does
not exploit. We give this finding the status it deserves: clean math
beats accidental signal only when the consumer is also rebuilt to read
the cleaner signal.
\end{abstract}

\tableofcontents

\section{Introduction}

E~prover's saturation loop is parametrised by a heuristic that assigns each
unprocessed clause a real weight; a min-heap selects the next clause to
process. The choice of heuristic shapes which proofs are found within a
fixed time budget. AI4REASON's \texttt{mzr02} strategy (Jakub\r{u}v and
Urban, IJCAR~2020) is a Mizar-tuned hand-engineered heuristic which on our
benchmark slice (\lpath{chainy_val_clean}, 800 problems from
the extended MPTP~5k validation set) solves $272/800$ in $30$~s per problem.
This is our single-heuristic baseline.

The Heuristic Control Block (HCB) mechanism in E permits round-robin
\emph{interleaving} of multiple weight functions inside the same prover
process: every clause is scored by every heuristic, and the scheduler
rotates through them with configurable multipliers. Crucially, all
heuristics share the same unprocessed-clause set, so a selection made by
heuristic $A$ generates child clauses that heuristic $B$ may later select.
In this paper, ``$+17$ over \texttt{mzr02}'' will always mean: an
interleaved configuration which mixes \texttt{mzr02}'s 9 entries (multiplier
sum $9$) with one ABGC entry at multiplier $3$, producing a $9{:}3$ rotation
($25\,\%$ ABGC clause picks). The score is a property of the \emph{interleaved}
process, not of either heuristic in isolation. Solo ABGC scores far below
either baseline (Sec.~\ref{sec:act1-solo}).

We address two questions:
\begin{itemize}
\item \textbf{Does ABGC, as interleaved with \texttt{mzr02}, beat \texttt{mzr02}-alone?}
      Yes, by $+17$ problems ($+6.2\,\%$ relative).
\item \textbf{Does a principled Hamilton--Jacobi--Bellman / Navier--Stokes
      reformulation of ABGC further improve the result?}
      No, in this paper's empirical regime: the principled transport
      loses the lift. We characterize the cause as a signal/consumer mismatch.
\end{itemize}

Both findings are reproducible from the artifacts listed in
Sec.~\ref{sec:reproducibility}.

\paragraph{Sibling paper.}
This paper is the saturation-loop sibling of our premise-selection paper
\emph{PLN as a Common Substrate for Naive Bayes and k-NN Premise Selection}
(\lpath{papers/PLN-kNN-NB.tex}), which addresses the upstream question of
\emph{which premises} to feed E on the same MPTP~5k slice. The two scopes
do not overlap: PLN-kNN-NB chooses the axioms; ABGC chooses, within
saturation, which clauses to process next.

\paragraph{Layout.}
Sec.~\ref{sec:architecture} describes the ABGC controller as a 14-bit bucket
abstraction with a static H-term and a Hodge-projected flow. Sec.~\ref{sec:act1}
reports the headline empirical result: $289/800$ in the hybrid, with set-diff
analysis and a solo-ABGC ablation. Sec.~\ref{sec:act2} reports the
HJB--NS principled refactor: eight Sub-phase~A configurations sweeping the
solver, sign convention, source budget, and pressure signal; three Sub-phase~B
configurations wiring viscosity, advection, and diffusion; and the SEGV
diagnosis. Sec.~\ref{sec:discussion} characterises why the cleaner controller
underperforms the accident: the consumer (the H-term) was implicitly tuned to a
signal regime that the principled controller removes by construction.
Sec.~\ref{sec:related} situates the work alongside MaLARea / MaSh / ENIGMA and
the watchlist tradition. Sec.~\ref{sec:conclusion} states what is open.
Sec.~\ref{sec:reproducibility} gives the heuristic string, the runner
invocation, the CSV paths backing each reported number, and the build hash.

\section{The ABGC architecture}\label{sec:architecture}

ABGC is implemented as a single E weight function, \lid{AdvWeight}, in
\lpath{HEURISTICS/che_abgc.c} (the source file referenced throughout, with
header \lpath{HEURISTICS/che_abgc.h}). The controller has four parts:
a bucket abstraction over clauses, a per-bucket dye field, a directed
parent-to-child flow on the bucket graph, and an H-term that combines static
clause features with three flow-derived signals.

\subsection{Bucket signatures}

Every clause is projected onto a 14-bit signature; this gives an address
space of $2^{14} = 16384$ logical buckets, but the controller allocates at
most \lid{ABGC_BMAX} $= 1024$ active bucket slots, hashed from the
14-bit signature with an overflow slot at (\lid{ABGC_BMAX} $-\, 1$). The
signature packs the following features into bit fields of width
$\{1,1,1,3,3,3,1,1\}$:
\begin{itemize}
\item is-goal (1 bit) --- did this clause come from a negated conjecture path?
\item is-unit (1 bit) --- is the clause a single literal?
\item is-ground (1 bit) --- are all terms ground?
\item depth bin (3 bits) --- 8 bins of clause proof-depth
\item size bin (3 bits) --- 8 bins of clause weight (log-spaced)
\item conjecture-overlap bin (3 bits) --- 8 bins of symbol-overlap with the
      conjecture's relevant function-symbol set
\item polarity-shape lo / hi (2 bits) --- coarse shape of the positive/negative
      literal distribution.
\end{itemize}
The hash from the 14-bit signature to a \lid{ABGC_BMAX} slot uses
Fowler--Noll--Vo (FNV-1a) modulo (\lid{ABGC_BMAX} $-\, 1$); collisions
land in the overflow slot. This gives the controller a bounded
state-vector (at most $\sim 1024$ active buckets) regardless of the
number of clauses E touches per problem.

\subsection{Per-bucket dye fields}

Each active bucket $b$ carries two non-negative scalars,
\begin{equation*}
   a_\mathrm{fwd}(b),\quad a_\mathrm{bwd}(b) \in [0, \infty),
\end{equation*}
called the forward and backward dyes. \emph{Forward} dye is injected at
bucket $b$ when an inserted clause is typed as Axiom or Hypothesis;
\emph{backward} dye is injected when the clause is typed as Conjecture or
NegConjecture. Both dyes decay multiplicatively at every batch boundary:
\begin{equation}
   a_{\bullet}(b) \;\mapsto\; (1 - \gamma_a\,\Delta t)\,a_{\bullet}(b),
\end{equation}
with $\gamma_a = 0.05$ and $\Delta t = 1.0$ by default. A bucket with both
dyes appreciably positive is a \emph{bridge bucket} --- a place where
axiom-rooted and goal-rooted derivations meet.

\subsection{Bucket-to-bucket flow}

For every clause inserted into the unprocessed set, the controller examines
the clause's derivation; if the derivation has a recorded first parent, the
parent clause's bucket $p$ and the child clause's bucket $b$ are linked by
incrementing $u_{p,b}$, a sparse flow matrix of dimension
\lid{ABGC_BMAX} $\times$ \lid{ABGC_BMAX}. The flow decays multiplicatively each batch with
$\gamma_u = 0.08$.

The net flow is the antisymmetric part
\begin{equation}
   F_{ij} = u_{ij} - u_{ji}.
\end{equation}
At each batch boundary the controller applies a discrete Hodge projection
to $F$ on the active-bucket subgraph (Sec.~\ref{sec:hodge}). The projected
$F$ is redistributed into a corrected $u$ which is then used by the H-term.

\subsection{The H-term}

Per-clause weight is the standard E ranking: lower weight wins. ABGC reads
off a real $H(\textrm{clause})$ which the caller subtracts from a base weight.
$H$ is a static linear combination of seven features:
\begin{align}
   H &\;=\; w_\mathrm{conj}\,h_\mathrm{conj}
            + w_\mathrm{depth}\,h_\mathrm{depth}
            + w_\mathrm{size}\,h_\mathrm{size}
            + w_\mathrm{goal}\,h_\mathrm{goal} \notag \\
     &\;\;\;\;{}+ w_\mathrm{meet}\,h_\mathrm{meet}
            + w_\mathrm{align}\,h_\mathrm{align}
            + w_\mathrm{pressure}\,h_\mathrm{pressure},
\end{align}
where
\begin{align}
   h_\mathrm{conj}     &= \texttt{conjecture\_overlap\_count}(\textrm{clause}) \\
   h_\mathrm{depth}    &= -\,\texttt{proof\_depth}(\textrm{clause}) \\
   h_\mathrm{size}     &= -\,\log_2\bigl(1 + \texttt{weight}(\textrm{clause})\bigr) \\
   h_\mathrm{goal}     &= [\,\texttt{is\_goal}(\textrm{clause})\,] \\
   h_\mathrm{meet}     &= \log\!\bigl(\varepsilon + a_\mathrm{fwd}(b) \cdot a_\mathrm{bwd}(b)\bigr) \\
   h_\mathrm{align}    &= u_{p \to b}\quad (p = \texttt{parent\_bucket}(\textrm{clause})) \\
   h_\mathrm{pressure} &= -\,p_\mathrm{Leray}(b),
\end{align}
with the Leray pressure $p_\mathrm{Leray}(b)$ produced as the dual variable
of the Hodge projection (Sec.~\ref{sec:hodge}). The weights
$(w_\mathrm{conj}, w_\mathrm{depth}, w_\mathrm{size}, w_\mathrm{goal},
 w_\mathrm{meet}, w_\mathrm{align}, w_\mathrm{pressure})$ are supplied
through E's heuristic configuration string; the configuration used for
all results in this paper is given in Sec.~\ref{sec:reproducibility}.

\subsection{Discrete Hodge projection}\label{sec:hodge}

We solve a Poisson equation on the active-bucket subgraph. Let
$L = D - A$ be the graph Laplacian with $A_{ij} = 1$ if either
$u_{ij}$ or $u_{ji}$ is non-zero and $D = \mathrm{diag}(\deg)$.
For each batch, define
\begin{equation}
   \mathrm{div}(F)[i] \;=\; \sum_{j > i:\, A_{ij} = 1} F_{ij}.
\end{equation}
We solve
\begin{equation}
   L\,p \;=\; \mathrm{div}(F)
\end{equation}
in the zero-mean subspace on the \emph{active} sub-graph (vertices with
$\deg \geq 1$); isolated vertices are excluded from the system and their
$p$ component is fixed at zero. The Phase~4 hybrid uses Jacobi iteration with
a cap of $50$ iterations; the HJB--NS refactor uses conjugate gradient on the
same null-space-corrected system. The corrected flow is
\begin{equation}
   F^{\mathrm{sol}}_{ij} \;=\; F_{ij} \;-\; (p_j - p_i),
   \quad
   u^{\mathrm{(new)}}_{ij}
   \;=\;
   \tfrac12 (u_{ij} + u_{ji}) + \tfrac12 F^{\mathrm{sol}}_{ij},
   \;\;
   u^{\mathrm{(new)}}_{ji}
   \;=\;
   \tfrac12 (u_{ij} + u_{ji}) - \tfrac12 F^{\mathrm{sol}}_{ij}.
\end{equation}
The sign convention $-\!(p_j - p_i)$ \emph{adds} the divergent component of
$F$ rather than removing it; the Helmholtz-correct projection would use
$+\!(p_j - p_i)$. This sign choice is mathematically wrong but empirically
useful (Sec.~\ref{sec:discussion}); the alternative is examined in
Sec.~\ref{sec:act2}.

\subsection{HCB interleaving with \texttt{mzr02}}

Outside the AdvWeight function, the heuristic string passed to E is a
$9{:}3$ multiplier sum
\begin{center}
\lid{(9 mzr02 entries) + 3*AdvWeight(...)}
\end{center}
which causes E's round-robin scheduler to pick from \texttt{mzr02}'s
weight functions 9 out of every 12 selections and from \lid{AdvWeight} 3
out of every 12. The unprocessed-clause set is shared across all entries.
Bucket-state and Hodge-projected flow are recomputed lazily at every
\lid{batch_period} $= 128$ insertions, so the cost amortizes to roughly a
single bucket-graph Hodge solve per $\sim 100$ clause-level operations.

\section{Act 1 --- empirical achievement}\label{sec:act1}

\subsection{Dataset, prover, hardware}\label{sec:setup}

The benchmark slice is \lpath{chainy_val_clean} from the
extended MPTP~5k dataset (originally derived from MML 5.94.1493 via the
chainy filter; we use the same 800-problem slice as the PLN-kNN-NB
sibling paper for direct comparability). Per-problem CPU limit is $30$~s,
imposed both at the runner level (\lid{timeout 70s} hard wall) and inside
E (\lid{--cpu-limit=30}). All runs use four parallel workers
(\lid{parallel --jobs 4}). Each worker has a $10$~GiB virtual-memory cap
(\lid{ulimit -v 10485760}). E's binary is built from
\lpath{/home/zar/claude/eprover/} with the ABGC source files linked into
the heuristics library; the relevant flags and build steps are recorded in
Sec.~\ref{sec:reproducibility}.

\subsection{Baselines}\label{sec:baselines}

\begin{table}[t]
\centering
\begin{tabular}{lrr}
\toprule
Configuration & Solved & Notes \\
\midrule
E \lid{--auto-schedule} (portfolio) & 334 / 800 & strategy ceiling, not a single-heuristic baseline \\
\texttt{mzr02} alone & 272 / 800 & single Mizar-tuned heuristic, our reference baseline \\
\bottomrule
\end{tabular}
\caption{Baselines on \lpath{chainy_val_clean}, 30~s CPU, 4 jobs.
The portfolio number \lid{--auto-schedule} is included so the reader does
not mistake $272$ as the prover's empirical ceiling; portfolio
schedulers route across many heuristics and are not the comparator for a
single-heuristic improvement claim.}
\label{tab:baselines}
\end{table}

\subsection{Hybrid result}\label{sec:hybrid}

\begin{table}[t]
\centering
\begin{tabular}{lrrr}
\toprule
Configuration & Solved & vs.\ \texttt{mzr02}-alone & vs.\ portfolio \\
\midrule
\texttt{mzr02} alone & 272 / 800 & --- & $-62$ \\
\textbf{\texttt{mzr02} + ABGC ($9{:}3$ HCB hybrid)} & \textbf{289 / 800} & \textbf{+17} & $-45$ \\
\midrule
$+\,6.2\,\%$ rel.\ over \texttt{mzr02} & & $+ 17$ & \\
\bottomrule
\end{tabular}
\caption{Single-heuristic baseline vs.\ the ABGC--\texttt{mzr02} hybrid on
the same 800-problem slice. CSV: \lpath{atps/logs/abgc_mzr02_hybrid_chainy_val_30s.csv}
($289$); \lpath{atps/logs/baseline_mzr02_chainy_val_30s.csv} ($272$).}
\label{tab:hybrid}
\end{table}

The ABGC--\texttt{mzr02} hybrid solves $289/800$, lifting the single-heuristic
baseline by $+17$ problems ($+6.2\,\%$ relative). Reproducibility is via the
runner script invocation in Sec.~\ref{sec:reproducibility}; the CSVs cited in
Table~\ref{tab:hybrid} exist on disk and a per-problem solved-count
\lid{awk} pass reproduces the headline number exactly.

\subsection{Set-diff: where the $+17$ comes from}

The $+17$ is a \emph{net}: the hybrid solves some problems that
\texttt{mzr02}-alone does not, and vice-versa. Per-problem set-diff on the
two CSVs gives:
\begin{itemize}
\item hybrid-unique wins: $35$ problems
\item \texttt{mzr02}-unique wins: $18$ problems
\item shared wins: $254$ problems.
\end{itemize}
Net: $35 - 18 = +17$.

The hybrid-unique wins are not concentrated in any narrow MPTP family; they
spread across \lid{compts}, \lid{connsp}, \lid{funct}, \lid{relset}, and
\lid{zfmisc} prefixes. The \texttt{mzr02}-unique losses are similarly
spread. The lift is a genuine breadth effect, not a narrow win on a
particular kind of problem.

\subsection{Solo-ABGC ablation}\label{sec:act1-solo}

Removing \texttt{mzr02} from the heuristic string and running ABGC
alone collapses the score. Solo-ABGC peaks in our exploration at roughly
$200$/$800$, well below either \texttt{mzr02}-alone ($272$) or the hybrid
($289$). This is not a failure of the controller; it is a feature of
the design. ABGC's dye field is injected only at axiom and conjecture
clauses, so derived clauses carry dye only through bucket-level transfer
(parent$\to$child flow), which is a sparse and noisy signal. In a hybrid,
\texttt{mzr02} provides the dense exploration; ABGC's role is to bias
toward bridges --- buckets where forward and backward dyes meet --- which
mzr02 by itself does not target. The HCB rotation lets ABGC's
``unusual'' picks generate children that \texttt{mzr02} then selects three
rounds later. The $+17$ is an interleaving effect, not a solo property.

\section{Act 2 --- HJB--NS principled refactor}\label{sec:act2}

The Phase~4 hybrid achieves its $+17$ with a controller that is, by the
standards of the underlying PDE skeleton (Goertzel \& Stay's HJB--NS
tutorial), \emph{ad hoc}: dye injection is unbounded (no source cap),
the projection's Jacobi iteration is run with a 50-iteration ceiling that is
often hit at residual $\approx 3$, the projection sign convention is
mathematically backwards, and there is no real transport step --- the
controller only \emph{decays} $u$ between batches.

A principled refactor of the controller along the Navier--Stokes /
HJB--NS skeleton suggests the following changes (per Sec.~\ref{sec:setup}
of the HJB--NS tutorial, Goertzel \& Stay):
\begin{align}
   \partial_t u + (u\!\cdot\!\nabla)u - \nu \Delta u
       &= -\nabla p - f, \quad \nabla\!\cdot\!u = 0 \\
   \partial_t a + \nabla\!\cdot\!(au)
       &= D \Delta a + s - \gamma a.
\end{align}
Mapped onto the bucket controller, this prescribes:
\begin{enumerate}
\item replace Jacobi by conjugate gradient (CG) so the projection actually converges;
\item enforce the Helmholtz-correct sign $F_\mathrm{sol} = F + (p_j - p_i)$;
\item add a per-batch source cap so injection is bounded;
\item add Laplacian viscosity $-\nu \Delta u$ on the net-flow field;
\item add mass-conserving upwind advection $\nabla\!\cdot\!(au)$ for the dyes;
\item add graph-Laplacian diffusion $D \Delta a$ on the dyes.
\end{enumerate}
We separate these into two sub-phases. Sub-phase~A keeps the controller
projection-only but sweeps the four solver / signal axes (CG vs Jacobi,
sign, source cap, pressure signal). Sub-phase~B wires the three transport
steps. Both were implemented and run end-to-end on the same 800-problem
slice.

\subsection{Sub-phase~A --- solver and signal sweeps}\label{sec:phaseA}

Five distinct configurations were evaluated on the full slice. Each
configuration is one single-variable swap relative to the Phase~4 hybrid
baseline.

\begin{table}[t]
\centering
\small
\begin{tabular}{llllll}
\toprule
Run & Solver & Sign & Sources & $h_\mathrm{pressure}$ & Solved \\
\midrule
Phase 4 hybrid & Jacobi & (b) $F-(p_j{-}p_i)$ & unbounded & Leray-$p$ & 289 \\
A-v5 (sign revert) & CG & (b) & bounded ($16$) & Leray-$p$ & 252 \\
A-v6 (cap removed) & CG & (b) & uncapped & Leray-$p$ & 253 \\
A-v7 (EMA) & CG & (b) & uncapped & EMA & 246 \\
A-v8 (Jacobi probe) & Jacobi & (b) & uncapped & EMA & 245 \\
\bottomrule
\end{tabular}
\caption{Sub-phase~A configurations on \lpath{chainy_val_clean}.
``Sign'' refers to the projection update:
(a) is the Helmholtz-correct $F_\mathrm{sol} = F + (p_j - p_i)$;
(b) is the divergence-amplifying $F_\mathrm{sol} = F - (p_j - p_i)$.
``EMA'' is the bucket-population-rate exponential moving average, an
alternative to Leray pressure for $h_\mathrm{pressure}$.
CSVs (all under \lpath{atps/logs/}):
\lpath{abgc_phaseA_v5_signrevert_chainy_val_30s.csv} (252);
\lpath{abgc_phaseA_v6_uncapped_chainy_val_30s.csv} (253);
\lpath{abgc_phaseA_v7_realema_chainy_val_30s.csv} (246);
\lpath{abgc_phaseA_v8_jacobi_chainy_val_30s.csv} (245).}
\label{tab:phaseA}
\end{table}

Reading off the deltas (Table~\ref{tab:phaseA}):
\begin{itemize}
\item \textbf{Sign convention} (A-v4 vs A-v5; not all rows shown):
      Helmholtz-correct sign $\to$ $-10$ problems.
\item \textbf{Source cap} (A-v5 vs A-v6): removing the cap of $16$ adds $+1$.
\item \textbf{Pressure signal} (A-v6 vs A-v7): swapping Leray-$p$ for
      bucket-population EMA costs $7$ problems.
\item \textbf{Solver} (A-v7 vs A-v8): swapping CG for Jacobi (at fixed
      everything-else) is essentially neutral: $-1$ problem. The solver
      is not load-bearing.
\end{itemize}
A-v6 is the best Sub-phase~A configuration at $253/800$; even with every
projection invariant honestly maintained (CG, isolated-node handling,
zero-mean re-projection), the score sits $36$ problems below the
Phase~4 hybrid. The math is correct; the controller has lost information.

\subsection{Sub-phase~B --- transport wired}\label{sec:phaseB}

Sub-phase~B wires the three transport steps onto the A-v6 baseline. The
batch update becomes
\begin{enumerate}
\item decay $a$, $u$;
\item build edge-list (lazily, from current support graph);
\item Laplacian viscosity on $F = u - u^{\top}$;
\item Hodge projection (CG);
\item mass-conserving upwind advection of $a_\mathrm{fwd}$ and $a_\mathrm{bwd}$;
\item graph-Laplacian diffusion of $a_\mathrm{fwd}$ and $a_\mathrm{bwd}$.
\end{enumerate}
Defaults: $\nu = 0.1$, $D = 0.04$, CFL safety factor $0.5$.

\paragraph{Sub-phase B v1 --- SEGV root-cause.}
The first wired run did not complete: $209/800$ problems returned
SZS status \lid{Unknown} after fewer than $1$~second of wall time, four
problems hung for $5+$ hours each, and the headline score was $164/800$.
A core dump from \lid{coredumpctl} placed the crash inside
\lid{ClauseSetExtractEntry} in E's clause-heap machinery
(\lpath{CLAUSES/ccl_clausesets.c}), not inside ABGC --- which is the
classic signature of a downstream consumer dereferencing a corrupted
comparator. The root cause was traced to
\lid{adv_advect_dyes}'s CFL clamp: the per-edge CFL bound is correct for
\emph{one} outflow edge, but a bucket $i$ with $k$ outflow edges loses
dye from \emph{every} edge, and the per-edge bound does not enforce the
per-node total. With $\sim 10$ outflow edges from a high-activity bucket,
the total outflow exceeds $a_\mathrm{fwd}(i)$, driving $a_\mathrm{fwd}(i)$
negative. Then $h_\mathrm{meet} = w_\mathrm{meet}\cdot\log(\varepsilon
+ a_\mathrm{fwd}\cdot a_\mathrm{bwd})$ produces $\log(\textrm{negative}) =
\mathrm{NaN}$; the NaN weight is inserted into E's clause heap; the heap
ordering is corrupted (NaN compares falsely against every other double);
\lid{ClauseSetExtractEntry} eventually dereferences a stale pointer
through the corrupted heap and SIGSEGVs.

\paragraph{Sub-phase B v2 --- fix.}
Three independent guards were added.
\begin{itemize}
\item In \lid{adv_advect_dyes}: replace the per-edge CFL clamp with a
      \emph{per-node-total} proportional share. Each outflow edge $e = (i,j)$
      transports
      $\phi_e = \mathrm{cfl\_safety}\cdot(\Delta t\, u_{ij}\,/\,\!\sum_{e' \in\, \mathrm{out}(i)} \Delta t\, u_{ie'})\cdot a_\mathrm{fwd}(i)$,
      so that $\sum_e \phi_e = \mathrm{cfl\_safety}\cdot a_\mathrm{fwd}(i) \leq a_\mathrm{fwd}(i)$
      by construction.
\item In \lid{adv_diffuse_dyes}: scale the diffusion coefficient per-edge by
      the larger endpoint degree, so the explicit-Euler stability bound
      $D \cdot \max(\deg) \leq 1/2$ holds even on high-degree nodes.
\item In \lid{adv_compute_H}: clamp $a_\mathrm{fwd}$ and $a_\mathrm{bwd}$ to
      $[0, \infty)$ before $\log(\,\cdot\,)$, and finite-check the final
      $H$ before returning. This is a last-line-of-defense guard against any
      future code path producing a non-comparable weight.
\end{itemize}
After the fix, zero \lid{Unknown} statuses are observed on the full slice.

\paragraph{Sub-phase B v2 / v3 results.}
\begin{table}[t]
\centering
\begin{tabular}{lllll}
\toprule
Run & Sign & Transport & $|u|_\infty$ behaviour & Solved \\
\midrule
A-v6 (reference) & (b) & none & $\gamma_u$-bounded & 253 \\
B-v2 & (b) & viscosity + advect + diffuse & overflow $\to$ inf & 246 \\
B-v3 & (a) Helmholtz-correct & viscosity + advect + diffuse & bounded & 238 \\
\bottomrule
\end{tabular}
\caption{Sub-phase~B on the same slice. Sign (a) bounds $|u|_\infty$
mathematically (the projection is divergence-removing) but loses an
additional $\sim 8$ problems versus sign (b). Sign (b) carries the
divergent signal that the H-term reads, but with viscosity + advection
recycling flow into the projection, the per-batch amplification
compounds and $|u|_\infty$ overflows. The score in the (b)-overflow
case is still meaningful: $246$ is not a random number, it is the
floor of how badly the H-term can perform when the align term is
saturated. CSVs:
\lpath{abgc_phaseB_v2_chainy_val_30s.csv} (246);
\lpath{abgc_phaseB_v3_signa_chainy_val_30s.csv} (238).}
\label{tab:phaseB}
\end{table}

\subsection{Sub-phase B telemetry}\label{sec:telemetry}

Per-batch telemetry is logged to
\lpath{atps/logs/abgc_telemetry/phaseB_v2.log}. Key macro statistics across
the $776$ PIDs that produced telemetry:
\begin{itemize}
\item \emph{Velocity-field magnitude}: median \lid{max_u_l1} reaches
      $9.7\times 10^{275}$; mean is $\infty$. On $99.2\,\%$ of PIDs,
      \lid{max_u_l1} eventually exceeds $1$. Sign~(b) plus clean CG
      convergence plus transport's flow recycling produces a true
      finite-time blowup at the double-precision floating ceiling.
\item \emph{Mixed-dye bridges}: median \lid{max_bridge_overlap} $= 3556$,
      max $= 88{,}000$. In Sub-phase~A configurations the same statistic
      was effectively zero, because $a_\mathrm{fwd}$ and $a_\mathrm{bwd}$
      lived in disjoint axiom-rooted and goal-rooted bucket sets with no
      transport to mix them. So the transport ``works'': bridges do form.
      The advection step is doing what it claims; the meet term has
      a non-trivial signal to read.
\item \emph{Projection convergence}: mean projection iterations
      $= 48.3 / 50$, mean residual $\approx 245$ at termination. CG is
      hitting the iteration cap on most batches, because the
      transport step on the previous batch produced fresh divergence
      that the current projection has to absorb. The
      under-convergence by itself is not catastrophic --- E's saturation
      loop is robust to noisy weights --- but together with the
      sign-(b) amplification it produces the overflow above.
\item \emph{Progress scalar}: median \lid{final_progress_ema} $= -0.21$
      (mostly negative); only $21.9\,\%$ of PIDs reach
      \lid{progress_ema} $> 0.05$. The controller's own progress signal
      reports decline, consistent with the score loss.
\end{itemize}

\section{Discussion --- why accidental signals beat clean math here}\label{sec:discussion}

\subsection{The divergence-amplifying sign convention is information}

In the discrete Hodge projection, the divergence-free part of a
co-chain $F$ on the bucket graph is
$F_\mathrm{sol}^{(a)} = F - \nabla p$. The correct sign convention
under the code's representation of $L$ and $\nabla$ is (a),
$F_\mathrm{sol}^{(a)} = F + (p_j - p_i)$ on the
indexed edges (Sec.~\ref{sec:hodge}). The Phase~4 hybrid uses sign~(b),
$F_\mathrm{sol}^{(b)} = F - (p_j - p_i)$.

Under the substitution $p_\mathrm{code}\to -p_\mathrm{true}$ that swaps
the two signs, one can see that
$F_\mathrm{sol}^{(b)} = F + \nabla p_\mathrm{true}$, i.e., the
divergent part of $F$ is \emph{added}, not subtracted. The divergent part
of $F$ on the bucket graph is precisely
$\nabla\!\cdot\!u\,|_b = \sum_j F_{bj}$: the net rate at which clauses
are leaving bucket $b$ across all parent-child links. \emph{Clause
production rate by bucket} is, intuitively, exactly what a clause-selection
heuristic should reward when picking the next clause from a producing
bucket: ``this bucket is making children fast; pick more of its parents.''
The sign-(b) convention writes this rate into $u_{p\to b}$ and the
H-term reads it through $h_\mathrm{align}$.

Sign-(a) removes the divergent part of $F$ exactly. The post-projection
$u_{p\to b}$ then carries only the divergence-free residual --- loops and
cycles on the bucket graph --- which is much weaker as a clause-production
signal.

\subsection{Under-convergence as the second piece of accidental information}

The Phase~4 hybrid runs the Jacobi solver with an iteration cap of $50$
and accepts the iterate as-is even when the residual is still
$\approx 3$ in the L${}^2$ norm of \lid{div(F)}. At active-graph sizes
of $200$--$400$ buckets, Jacobi on the graph Laplacian does not converge
under the iteration cap; the controller's $p$ is therefore
\emph{partial} --- close to the true potential on small-$|p|$ regions,
far from it on high-divergence regions. The sign-(b) update
$F_\mathrm{sol}^{(b)} = F - (p_j - p_i)$ with a small-magnitude partial
$p$ leaves most of $F$'s divergent component intact and only locally
attenuates the spikes.

This is a happy accident: Jacobi as a poor solver acts as a
\emph{frequency filter} on the divergence signal, smoothing high-frequency
modes but preserving the low-frequency clause-production rate that the
H-term wants. When we replace Jacobi with CG (Sub-phase~A) we get full
convergence in $\sim 10$ iterations to residual $\sim 10^{-6}$; the
sign-(b) update then writes the \emph{full} divergent signal, including
high-frequency artifacts. The H-term reads more loudly but the signal is
noisier; in the wash, we lose $\sim 10$ problems.

Adding transport (Sub-phase~B) closes the loop: viscosity smooths $u$
toward an equilibrium, advection redistributes dyes along $u$, and the
projection then has to absorb the freshly-produced divergence from the
redistribution. With sign~(b) this becomes a per-batch amplification
loop and $|u|_\infty$ blows up. With sign~(a) it stays bounded but
divergence-free, and the divergent signal the H-term wants is gone.
Either way, the H-term reads less.

\subsection{The H-term was implicitly tuned to the Phase~4 regime}

The conclusion we draw is structural, not numerical: the H-term shape
$\mathrm{meet} + \mathrm{align} + \mathrm{pressure}$ was tuned (by hand,
on small datasets, over many evenings) against the Phase~4 hybrid's
specific signal regime --- under-converged Jacobi projection, sign-(b)
divergence amplification, unbounded source injection. When we replace any
of these with its principled counterpart the H-term reads a structurally
different signal and the tuning becomes stale. A principled controller
delivers a principled signal; the consumer has to be re-tuned, perhaps
even re-designed, to read it. The cost of \emph{not} re-tuning is the
$+17$ lift.

\subsection{What kind of negative result is this?}

This is not a refutation of the HJB--NS framing for inference control.
The transport step does what its specification says it does: mass is
conserved, bridges form, projection convergence is honest. What it does
\emph{not} do is improve a heuristic whose H-term was implicitly fitted
to a different signal. The principled refactor exposes a hidden
brittleness in the original hybrid: the $+17$ depends sensitively on
the specific accident of sign~(b) plus under-convergence. A reviewer who
treats this controller as ``mzr02 with a clever Hodge projection''
underestimates how much hand-tuning the H-term encodes against the
projection's idiosyncrasies.

\subsection{What would a real HJB--NS clause-selection look like?}

We do not believe the right next step is to fiddle with the H-term to
recover the $+17$ on top of clean transport; that would just be
re-tuning to a new accident. The right next steps, in our view, are:
\begin{itemize}
\item \emph{Replace the static H-term with a learned reader} of the
      projected $u$ and the dye fields, with the divergence signal
      explicitly factored out as a separate ranking feature.
\item \emph{Move from bucket-level to per-clause Hodge projection on the
      derivation graph} (the actual DAG of generated clauses), so the
      Helmholtz decomposition operates on the object the H-term actually
      reads.
\item \emph{Replace the hand-tuned $\phi$-basis} ($h_\mathrm{conj}$,
      $h_\mathrm{size}$, etc.) with an ENIGMA-style learned $\phi$, so
      the consumer can adapt to whatever signal the principled transport
      provides.
\end{itemize}
None of these is done in this paper. The honest contribution of Act 2 is
the diagnosis, not the fix.

\section{Related work}\label{sec:related}

\paragraph{AI4REASON, MaLARea, MaSh, ENIGMA, MaSh-style.}
The single-heuristic baseline \texttt{mzr02} is from the AI4REASON
group's IJCAR~2020 paper on hand-tuned Mizar strategies
(Jakub\r{u}v and Urban, IJCAR 2020). MaLARea and MaSh are premise
selectors and are upstream of the saturation step we operate in.
ENIGMA learns clause-selection heuristics by gradient descent on a
$\phi$-basis; our static H-term is a hand-engineered analog. Our
sibling paper \emph{PLN as a Common Substrate for Naive Bayes and k-NN
Premise Selection} (\lpath{papers/PLN-kNN-NB.tex}) addresses the
premise-selection axis on the same 800-problem slice and is the most
direct point of comparison for the dataset and protocol.

\paragraph{Watchlist heuristics.}
Goal-directed watchlist heuristics in E and Vampire are the closest
prior art for ``bucket the unprocessed-clause set and bias toward
buckets that look productive.'' Our 14-bit bucket signature is more
coarse than a typical watchlist (which is per-symbol), and our
dye+flow mechanism makes the bias quantitative rather than binary.

\paragraph{Fluid analogies in proof search.}
Goertzel and Stay's HJB--NS tutorial in
\lpath{literature/Goertzel_Ben/benxiv/Fluid-AI/HJB-NS-tutorial.pdf} is the
direct motivation for the Sub-phase B transport; we have implemented it
as faithfully as the bucket-level setting allows and reported the
result honestly. The framing as such is sound; whether it pays off
empirically depends on whether the consumer of the projection is built
to read it.

\section{Conclusion and open work}\label{sec:conclusion}

The static-H Phase~4 hybrid is a usable clause-selection heuristic:
$+17$ over \texttt{mzr02} on the chainy MPTP slice, $+6.2\,\%$
relative, reproducible from the CSV artifacts on disk. The HJB--NS
principled refactor does not exceed it; we have characterised this as
a signal/consumer mismatch rather than a failure of the framing. The
following questions are open and we do not address them here:
\begin{itemize}
\item Can the H-term be rebuilt to read a divergence-free signal as
      productively as it reads the divergence-amplified one? (We do not
      think the answer is yes for the current static H-term shape; an
      ENIGMA-style learned H-term might.)
\item Does the framing pay off at a finer granularity? Per-clause Hodge
      projection on the derivation DAG would put the projection at the
      level the H-term actually reads, which we have not implemented.
\item Does the framing pay off on a different benchmark? We have only
      run on chainy MPTP; harder TPTP slices or Mizar Mathematical
      Library splits might exhibit a different signal regime.
\end{itemize}

\section{Reproducibility}\label{sec:reproducibility}

\subsection{Heuristic string}

The HCB string passed to \lid{eprover -H} for the headline hybrid result
is the following $9{:}3$ multiplier sum
(\texttt{mzr02} entries first, AdvWeight last):

\begin{lstlisting}
(
 1*ConjectureTermPrefixWeight(DeferSOS,1,3,0.1,5,0,0.1,1,4),
 1*ConjectureTermPrefixWeight(DeferSOS,1,3,0.5,100,0,0.2,0.2,4),
 1*Refinedweight(PreferWatchlist,4,300,4,4,0.7),
 1*RelevanceLevelWeight2(PreferProcessed,0,1,2,1,1,1,200,200,2.5,9999.9,9999.9),
 1*StaggeredWeight(DeferSOS,1),
 1*SymbolTypeweight(DeferSOS,18,7,-2,5,9999.9,2,1.5),
 2*Clauseweight(PreferWatchlist,20,9999,4),
 2*ConjectureSymbolWeight(DeferSOS,9999,20,50,-1,50,3,3,0.5),
 2*StaggeredWeight(DeferSOS,2),
 3*AdvWeight(ConstPrio,1.0,8.0,2.0,1.0,25.0,3.0,5.0,10.0)
)
\end{lstlisting}
The \lid{AdvWeight} arguments are
$(\alpha, w_\mathrm{conj}, w_\mathrm{depth}, w_\mathrm{size}, w_\mathrm{goal},
   w_\mathrm{meet}, w_\mathrm{align}, w_\mathrm{pressure}) =
(1.0, 8.0, 2.0, 1.0, 25.0, 3.0, 5.0, 10.0)$.

\subsection{Runner invocation}

All scores reported are produced by the runner script
\lpath{atps/scripts/run_baseline_prover.sh}, invoked as
\begin{lstlisting}
ABGC_LOG_PATH=atps/logs/abgc_telemetry/PHASE.log
  atps/scripts/run_baseline_prover.sh \
  PHASE_chainy_val_30s \
  'eprover --free-numbers --definitional-cnf=24 --split-aggressive
   --simul-paramod --forward-context-sr --destructive-er-aggressive
   --destructive-er --prefer-initial-clauses -tKBO6 -winvfreqrank -c1
   -Ginvfreq -F1 --delete-bad-limit=150000000
   -WSelectMaxLComplexAvoidPosPred -H"<<heuristic string above>>"
   --cpu-limit={TIMEOUT} {PROBLEM}' \
  atps/datasets/extended_mptp5k/chainy_val_clean \
  30 \
  atps/logs/PHASE_chainy_val_30s.csv \
  4
\end{lstlisting}
with \lpath{ulimit -v 10485760} (10~GiB virtual-memory cap) imposed at the
runner-wrapper level.

\subsection{CSV artifacts}

\begin{longtable}{L{0.55\linewidth} r r}
\toprule
CSV path (under \lpath{atps/logs/}) & Solved & Total \\
\midrule
\endhead
\lpath{baseline_e_auto_sched_chainy_val_30s.csv} & 334 & 800 \\
\lpath{baseline_mzr02_chainy_val_30s.csv} & 272 & 800 \\
\lpath{abgc_mzr02_hybrid_chainy_val_30s.csv} & 289 & 800 \\
\lpath{abgc_phaseA_v5_signrevert_chainy_val_30s.csv} & 252 & 800 \\
\lpath{abgc_phaseA_v6_uncapped_chainy_val_30s.csv} & 253 & 800 \\
\lpath{abgc_phaseA_v7_realema_chainy_val_30s.csv} & 246 & 800 \\
\lpath{abgc_phaseA_v8_jacobi_chainy_val_30s.csv} & 245 & 800 \\
\lpath{abgc_phaseB_v2_chainy_val_30s.csv} & 246 & 800 \\
\lpath{abgc_phaseB_v3_signa_chainy_val_30s.csv} & 238 & 800 \\
\bottomrule
\end{longtable}
Each row is verifiable as
\lid{awk -F, 'NR>1 \&\& \$2==1 \{n++\} END\{print n+0\}' <csv>}.

\subsection{Source tree}

\begin{itemize}
\item \lpath{HEURISTICS/che_abgc.c} (E~prover heuristics directory) ---
      the ABGC controller. The active \lid{adv_batch_update} calls only
      the Hodge projection by default; the Sub-phase B transport functions
      (\lid{adv_build_edge_list}, \lid{adv_viscosity_step},
      \lid{adv_advect_dyes}, \lid{adv_diffuse_dyes}) are present in the file
      and statically callable, but not wired in the current production state.
\item \lpath{HEURISTICS/che_abgc.h} --- type definitions.
\item \lpath{HEURISTICS/che_heuristics.c} --- the registration site
      \lid{HCBStdInit} where \lid{AdvWeightInit} is exposed to the
      heuristic-string parser.
\end{itemize}

\bibliographystyle{plain}
\begin{thebibliography}{10}

\bibitem{ai4reason-ijcar20}
Jan Jakub\r{u}v and Josef Urban.
\newblock Hammering Mizar by learning clause guidance (Mizar40 / mzr02).
\newblock In \emph{IJCAR 2020}.

\bibitem{e-prover}
Stephan Schulz.
\newblock System description: E 1.8.
\newblock In \emph{LPAR-19}, volume 8312 of \emph{LNCS}, pages 735--743.
        Springer, 2013.

\bibitem{goertzel-stay-hjbns}
Ben Goertzel and Mike Stay.
\newblock HJB--Navier--Stokes for inference control: a tutorial.
\newblock Working note, \lpath{literature/Goertzel_Ben/benxiv/Fluid-AI/HJB-NS-tutorial.pdf}, 2025.

\bibitem{mash}
Daniel K\"uhlwein and Josef Urban.
\newblock MaSh: Machine learning for Sledgehammer.
\newblock In \emph{ITP 2013}.

\bibitem{enigma}
Jan Jakub\r{u}v and Josef Urban.
\newblock ENIGMA: Efficient learning-based inference guiding machine.
\newblock In \emph{CICM 2017}.

\bibitem{mptp}
Josef Urban.
\newblock MPTP 0.2: Design, implementation, and initial experiments.
\newblock \emph{Journal of Automated Reasoning}, 37:21--43, 2006.

\bibitem{pln-knn-nb}
Mettapedia Working Notes.
\newblock PLN as a common substrate for Naive Bayes and k-NN premise selection.
\newblock In this directory: \lpath{papers/PLN-kNN-NB.tex}.

\bibitem{helmholtz-decomp-graphs}
Lek-Heng Lim.
\newblock Hodge Laplacians on graphs.
\newblock \emph{SIAM Review}, 62(3):685--715, 2020.

\end{thebibliography}

\end{document}
